\documentclass[12pt]{article}
\usepackage{amsmath}
\usepackage{amssymb}
\usepackage{amsthm}
\usepackage{cancel}
\usepackage{enumitem}
\usepackage{esdiff}
\usepackage{graphicx}
\usepackage[version=4,textfontname=sffamily,mathfontname=mathsf]{mhchem}
\usepackage{multicol}
\usepackage{siunitx}
\usepackage{ulem}
% \usepackage{pgfplots}
\usepackage{wrapfig}
\usepackage{xcolor}

\newcommand{\e}[1]{e^{#1}}
\newcommand{\ei}[1]{e^{i\left(#1\right)}}
\newcommand{\E}[1]{\times 10^{#1}}
\newcommand{\tagref}[1]{\tag{\ref{#1}}}

\renewcommand\qedsymbol{TENA}

\title{
    Homework \#14
    \\  \small
    PHYS 4D: Modern Physics
    \\
    Problems from Krane's textbook
    }
\author{Donald Aingworth IV}
\date{April 27, 2026}

\begin{document}
\maketitle
\DeclareSIUnit{\amu}{amu}
\DeclareSIUnit{\light}{c}
\DeclareSIUnit{\lightyear}{c\cdot yr}
\DeclareSIUnit{\year}{yr}
\DeclareSIUnit{\yr}{yr}

\section{Questions, Chapter 2}
    \subsection{Question 1}
        Explain in your own words what is meant by the term ``relativity''. 
        Are there different theories of relativity?

        \subsubsection{Answer}
            Relativity is the description of the observation of events with different positions and velocities.
            There are multiple theories of relativity, all of which postulate that physics works the same in all reference frames.
            Classical relativity posits no limit to speed and considers a single observation of the speed of time.
            Special relativity builds on classical relativity and posits thar the speed of light must be identical in all reference frames.
            General relativity also exists, and concerns gravity and the shape of spacetime.
        \pagebreak

    \subsection{Question 8}
        How does relativity combine space and time coordinates into spacetime?
        
        \subsubsection{Answer}
            It uses the pythagorean theorem, the speed of the moving observer, and an electron clock to relate the velocities.
            It uses equations, most of which involve $u^2/c^2$, to relate the time elapsed in one time frame, the velocity of another time frame relative to the first, and the time elapsed in the other time frame.

    \subsection{Question 10}
        Explain in your own words the terms \textit{time dilation} and \textit{length contraction}.

        \subsubsection{Answer}
            Time dilation refers to moving objects appearing to experience time slower than objects moving at lesser speeds.
            Length contraction refers to moving objects appearing to have significantly shorter length relative to their length at rest. 
            This is usually only observable in objects moving at relatively incredibly high speeds, usually close to the speed of light.
    \pagebreak

\section{Problem 6} % 2.4
    An astronaut must journey to a distant planet, which is 200 light-years from Earth. 
    What speed will be necessary if the astronaut wishes to age only 10 years during the round trip?

    \subsection{Solution}
        Take the formula for time dilation.
        \begin{equation}
            \Delta t    =   \frac{\Delta t_0}{\sqrt{1 - u^2/c^2}}
        \end{equation}

        Define time in years and velocity in terms of light-years.
        We can replace $\Delta t$ with $\frac{L}{u}$.
        Substitute in with the equation for relative $L$.
        Then, solve for $u$.
        I'll be using the quadratic formula and removing values of $c$ for the sake of cleanliness.
        \begin{gather}
            \Delta t    =   \frac{\Delta t_0}{\sqrt{1 - u^2/c^2}}\\
            \frac{L}{u}    =   \frac{\Delta t_0}{\sqrt{1 - u^2/c^2}}\\
            L_0 \sqrt{1 - u^2/c^2} * \sqrt{1 - u^2/c^2} = u\,\Delta t_0\\
            L_0 - L_0 \frac{u^2}{c^2} = u\,\Delta t_0\\
            0   =   u^2 + \frac{\Delta t_0}{L_0} u - 1
                =   u^2 + \frac{5}{200} u - 1\\
            \begin{align}
                u   &=  \frac{-\frac{1}{40} \pm \sqrt{\left( \frac{1}{40} \right)^2 + 4}}{2}c
                    =   \frac{-\frac{1}{40} \pm 2.000156}{2}c\\
                    &=  \frac{1.975}{2}c
                    =   \boxed{0.988c}
            \end{align}
        \end{gather}
    \pagebreak

\section{Problem 8}
    High-energy particles are observed in laboratories by photographing tracks they leave in certain detectors; the length of the track depends on the speed of the particle and its lifetime. 
    A particle moving at 0.995c leaves a track 1.25 mm long. 
    What is the proper lifetime of the particle?

    \subsection{Solution}
        First, find the observed lifetime of the particle.
        \begin{equation}
            \Delta t_o = \frac{\Delta x}{v}
                =   \frac{1.25\,\unit{\milli\meter}}{0.995 * 2.998\E{8}\,\unit{\meter/\second}}
                =   4.19\E{-12}\,\unit{\second}
        \end{equation}

        Put that into the formula for the proper lifetime of the particle.
        \begin{align}
            \Delta t    &=  \frac{\Delta t_o}{\sqrt{1 - u^2/c^2}}
                =   \frac{4.19\E{-12}\,\unit{\second}}{\sqrt{1 - (0.995)^2}}
                =   \frac{4.19\E{-12}\,\unit{\second}}{\sqrt{1 - 0.990025}}\\
                &=  \frac{4.19\E{-12}\,\unit{\second}}{\sqrt{0.009975}}
                =   \frac{4.19\E{-12}\,\unit{\second}}{0.09987}
                =   \boxed{4.195\E{-11}\,\unit{\second}}
        \end{align}
    \pagebreak

\section{Problem 9}
    Carry out the missing steps in the derivation of Eq. 2.17.
    \begin{equation}
        \label{eq:2.17} \tag{2.17}
        v = \frac{v' + u}{1 + v'u/c^2}
    \end{equation}

    \subsection{Solution}
        We begin, as the book says, with Equations (2.15) and (2.16) and solve them for $\Delta t_1$ and $\Delta t_2$.
        \begin{gather}
            \tag{2.15}
            v \Delta t_1 = L + u\Delta t_1\\
            (v - u) \Delta t_1 = L\\
            \Delta t_1 = \frac{L}{v - u}\\
            \tag{2.16}
            c \Delta t_2 = L - u \Delta t_2\\
            (c + u) \Delta t_2 = L\\
            \Delta t_2 = \frac{L}{c + u}
        \end{gather}

        Add these together to get the total time elapsed.
        \begin{align}
            \Delta t    &=  \Delta t_1 + \Delta t_2
                =   \frac{L}{v - u} + \frac{L}{c + u}
                =   L \frac{c + u + v - u}{(v - u)(c + u)}\\
                &=  L \frac{c + v}{(v - u)(c + u)}
        \end{align}

        Relate this to both $\Delta t_o$ and $L_o$ from prior formulas.
        \begin{gather}
            \Delta t = \frac{\Delta t_o}{\sqrt{1 - u^2/c^2}}\\
            L = L_o \sqrt{1 - u^2/c^2}\\
            \frac{\Delta t_o}{\sqrt{1 - u^2/c^2}} = L_o \sqrt{1 - u^2/c^2} \frac{c + v}{(v - u)(c + u)}\\
            \Delta t_o  =   L_o \frac{(c + v)(1 - u^2/c^2)}{(v - u)(c + u)}
        \end{gather}

        Relate this to the formula for the time interval passed relative to the clock itself, Eq. (\ref{eq:2.14}).
        \begin{gather}
            \label{eq:2.14} \tag{2.14}
            \Delta t_o = \frac{L_o}{v'} + \frac{L_o}{c}
                =   L_o \frac{v' + c}{v'c}\\
            L_o \frac{v' + c}{v'c} = L_o \frac{(c + v)(1 - u^2/c^2)}{(v - u)(c + u)}\\
            \frac{v' + c}{v'c (1 - u^2/c^2)} = \frac{c + v}{(v - u)(c + u)}
        \end{gather}

        \begin{gather}
            (v' + c)(v - u)(c + u) = v'c(c + v)\left( 1 - \frac{u^2}{c^2} \right)\\
            \begin{split}
                v'vc + v'vu - v'uc - v'u^2 &+ c^2v + cvu - c^2u - cu^2 \\
                    &= v'c^2 + v'vc - \frac{v'c^2u^2}{c^2} - \frac{v'vcu^2}{c^2}
            \end{split}
            \\
            v'vu + c^2v + cvu + \frac{v'vu^2}{c} = v'c^2 + v'uc + c^2u + cu^2 + v'u^2 - v'u^2 \\
            v\left( v'u + c^2 + cu + \frac{v'u^2}{c} \right) = v'c^2 + v'uc + c^2u + cu^2 \\
            v\left( c(c + u) + \frac{v'uc}{c} + \frac{v'u^2}{c} \right) = v'c(u + c) + cu(u + c) \\
            v\left( c(c + u) + \frac{v'u}{c}(c + u) \right) = c(u + c)(v' + u) \\
            v (c + u) \left( c + \frac{v'u}{c} \right) = c(u + c)(v' + u) \\
            v \left( 1 + \frac{v'u}{c^2} \right) = v' + u \\
            \boxed{v = \frac{v' + u}{1 + \frac{v'u}{c^2}}}
        \end{gather}
    \pagebreak

\section{Problem 10}
    Two spaceships approach the Earth from opposite directions.
    According to an observer on the Earth, ship A is moving at a speed of 0.753c and ship B at a speed of 0.851c. 
    What is the velocity of ship A as observed from ship B? 
    Of ship B as observed from ship A?
    
    \subsection{Solution}
        Recall the equation for classical velocity addition.
        \begin{equation}
            v = v' + u
        \end{equation}

        Start with the equation for relativistic velocity.
        \begin{equation}
            v = \frac{v' + u}{1 + \frac{v'u}{c^2}}
        \end{equation}

        Within the classical equation, suppose that $u = v_{BA}$, the velocity of B as perceived from A.
        We can call the velocity of A as perceived from earth $v_A$, a stand-in for $v'$.
        $v$ will then be replaced by the velocity of B as perceived from earth, or $v_B$.
        Plug these into the classical velocity addition equation and solve for $v_{BA}$.
        \begin{gather}
            v_B = v_A + v_{BA}\\
            v_{BA} = v_B - v_A
        \end{gather}

        Compare like terms with the classical relativity equation to find what replaces what when plugging into the spacial relativistic velocity equation.
        \begin{gather}
            v \to v_{BA};
            v' \to v_B;
            u \to -v_A\\
            \begin{align}
                v_{BA}  &=  \frac{v_B - v_A}{1 - \frac{v_A v_B}{c^2}}
                    =   \frac{0.851\,\unit{\light} - 0.753\,\unit{\light}}{1 - 0.851 * 0.753}
                    =   \frac{0.098}{0.359197}\,\unit{\light}
                    =   \boxed{0.273\,\unit{\light}}
            \end{align}
        \end{gather}

        For the speed of A as seen from B, take the negative, so $v_{AB} = \boxed{-0.273\,\unit{\light}}$.
    \pagebreak

\section{Problem 15} % 2.5 
    Observer $O$ fires a light beam in the $y$ direction ($v_y = c$).
    Use the Lorentz velocity transformation to find $v'_x$ and $v'_y$ and show that $O'$ also measures the value $c$ for the speed of light. 
    Assume that $O'$ moves relative to $O$ with velocity $u$ in the $x$ direction.

    \subsection{Solution}
        First define the velocity vector. \begin{equation}
            \vec{v} = \begin{pmatrix} v_x \\ v_y \\ v_z \end{pmatrix} = \begin{pmatrix} 0 \\ c \\ 0 \end{pmatrix}
        \end{equation}

        First find $v_x$.
        \begin{align}
            v_x'    &=  \frac{v_x - u}{1 - v_x u/c^2}
                =   \frac{0 - u}{1 - 0 u/c^2}
                =   -\frac{u}{1}
                =   -u
        \end{align}

        Then $v_y$.
        \begin{align}
            v_y'    &=  \frac{v_y \sqrt{1 - u^2/c^2}}{1 - v_x u/c^2}
                =   \frac{c \sqrt{1 - u^2/c^2}}{1 - 0 u/c^2}
                =   \sqrt{c^2 - u^2}
        \end{align}

        Lastly $v_z$ for good measure.
        \begin{equation}
            v_z'    =   \frac{v_z \sqrt{1 - u^2/c^2}}{1 - v_x u/c^2}
                =   \frac{0 \sqrt{1 - u^2/c^2}}{1 - 0 u/c^2}
                =   0
        \end{equation}

        To find the speed of light, invoke Pythagoras.
        \begin{equation}
            \left| \vec{v} \right|  =   \sqrt{(-u)^2 + \sqrt{c^2 - u^2}^2}\\
                =   \sqrt{u^2 + c^2 - u^2}
                =   \sqrt{c^2}
                =   c
        \end{equation}
        TENA
    \pagebreak

\section{Problem 18}
    A rod in the reference frame of observer $O$ makes an angle of $31\unit{\degree}$ with the $x$ axis. 
    According to observer $O'$, who is in motion in the $x$ direction with velocity $u$, the rod makes an angle of $46\unit{\degree}$ with the $x$ axis. 
    Find the velocity $u$.

    \subsection{Solution}
        If $O'$ is only moving along the $x$ axis, we don't even need to use the Lorentz transformation.
        Suppose the rod has a length $\ell$.
        In the $x$-axis from $O$, the rod extends a distance of $\ell\cos(31\unit\degree)$, which can be our value of $L_o$.
        In the $x$-axis from $O'$, the rod extends a distance of $\ell\cos(46\unit\degree)$, which can be our value of $L$.
        From here, we can use the length transformation equation to find $u$.
        \begin{gather}
            L = L_o \sqrt{1 - u^2/c^2}\\
            \frac{L^2}{L_o^2} = 1 - u^2/c^2\\
            \frac{u^2}{c^2} =   1 - \frac{L^2}{L_o^2}
                =   1 - \frac{\cos^2(46\unit\degree)}{\cos^2(31\unit\degree)}
                =   1 - \frac{0.48255}{0.73474}
                =   0.34323\\
            \boxed{u = 0.58586c}
        \end{gather}
\end{document}